CIRCULAR CONVOLUTION: FLIP + SHIFT

N-point circular convolution:

$$y[n]=\sum_{m=0}^{N-1}x[m]h[(n-m)\text{mod} N]$$

Indices wrap around mod N. Output always has N samples.

MAIN MEMORY TRICK

1. Keep x[m] FIXED, always in positions 0 through N-1.

2. CIRCULARLY FLIP h[m] to form h[(-m) mod N]. The h[0] entry stays at index 0.

3. For output n, shift the flipped h row RIGHT by n positions (wrapping around).

4. Multiply aligned entries by fixed x[m], then ADD to obtain y[n].

5. Repeat until n=N-1.

WORKED EXAMPLE: N=4

Fixed input:

$$x[m]=[1,2,3,4]$$

Original impulse response:

$$h[m]=[1,0,2,3]$$

CIRCULAR FLIP OF h

Original indices $m=0,1,2,3$:

$$h[m]=[h[0],h[1],h[2],h[3]]$$

Circularly flipped:

$$h[(-m)\text{mod}4]$$

$$=[h[0],h[3],h[2],h[1]]$$

$$=[1,3,2,0]$$

IMPORTANT: Flip keeps h[0] FIXED, then reverses the remaining positions. It is NOT ordinary reverse [3,2,0,1].

HOLD x FIXED FOR EVERY ROW

$$x=[1,2,3,4]$$

For n=0, shifted/flipped h:

$$h[(-m)\text{mod}4]=[1,3,2,0]$$

$$y[0]=1(1)+2(3)+3(2)+4(0)$$

$$y[0]=13$$

For n=1, RIGHT SHIFT flipped h by 1:

$$h[(1-m)\text{mod}4]=[0,1,3,2]$$

$$y[1]=1(0)+2(1)+3(3)+4(2)$$

$$y[1]=19$$

For n=2, RIGHT SHIFT flipped h by 2:

$$h[(2-m)\text{mod}4]=[2,0,1,3]$$

$$y[2]=1(2)+2(0)+3(1)+4(3)$$

$$y[2]=17$$

For n=3, RIGHT SHIFT flipped h by 3:

$$h[(3-m)\text{mod}4]=[3,2,0,1]$$

$$y[3]=1(3)+2(2)+3(0)+4(1)$$

$$y[3]=11$$

FINAL RESULT:

$$y[n]=[13,19,17,11]$$

CHECK: sum of circular output = (sum x)(sum h):

$$13+19+17+11=10(6)=60$$

THE FLIP / SHIFT / MULTIPLY TABLE

x always: [1, 2, 3, 4]

n=0 h row: [1, 3, 2, 0]; dot=13.

n=1 h row: [0, 1, 3, 2]; dot=19.

n=2 h row: [2, 0, 1, 3]; dot=17.

n=3 h row: [3, 2, 0, 1]; dot=11.

Each new row is the preceding h row shifted ONE position RIGHT with wraparound. Never move the x row.

DFT CIRCULAR CONVOLUTION THEOREM

$$X[k]=\text{DFT}_N(x[n])$$

$$H[k]=\text{DFT}_N(h[n])$$

$$y[n]=\text{IDFT}_N(X[k]H[k])$$

For linear convolution of finite signals of lengths Lx and Lh, zero-pad both to N at least:

$$N\ge L_x+L_h-1$$

Then N-point circular convolution equals the linear convolution with zero-padding. Without enough padding, circular convolution folds linear-convolution samples:

$$y_{circ}[n]=\sum_r y_{lin}[n+rN]$$

where the sum covers valid samples.
